\section{未来 Vignettes 与 Q\&A：从廉价智能到社会智能}

正式讲稿最后四张图把前面的系统设计外推到未来：能力成本下降、界面动态生成、教育与医疗个性化、故事创作关系改变。随后 22 分钟 Q\&A 给出更有价值的限制条件。讲义将愿景与回答合并阅读：每个“可以”后面都补上“还缺什么 evidence”。

\subsection{能力成本下降，但验证成本未必同步下降}

MMLU cost frontier 图表达的是历史趋势：在相似准确率下，推理价格快速下降。它不能证明成本永远指数下降，也不能说明 MMLU 代表真实工作。更重要的是，生成变便宜会增加输出数量；若验证仍依赖专家，系统总成本可能从 inference 转移到 review。

\lecturefigure{slide-57-mmlu-cost-frontier.jpg}{MMLU performance--cost frontier：同类能力的调用成本历史性下降}{00:45:46--00:47:40}

令任务成功率为 $q(c)$、单次推理成本为 $c$、人工验证成本为 $v$，单位可信成功成本可粗略写为
\[
C_{\text{trusted}}=\frac{c+v}{q(c)}.
\]
即使 $c$ 下降，若更复杂任务使 $v$ 上升或 $q$ 不稳定，可信结果并不会同速变便宜。产品应优化 $C_{\text{trusted}}$，而不是只报 token price。

\teachervoice{00:45:46--00:47:40，讲者说这张图已是“一年前做的”，MMLU 也不再是大家最关心的 benchmark；她用它表达 raw intelligence 变便宜。讲义提醒：这是 dated trend，不是无限外推定律。}

\subsection{Dynamic UI、个性化服务与故事关系}

上一节说明 raw intelligence 的调用成本下降，但低价输出只有进入合适界面才会产生用户价值。本节把成本趋势连接到三类 surface：随意图生成的动态 UI、承担高风险建议的教育/医疗入口，以及保留作者 agency 的故事协作。Dynamic generative UI 的愿景是界面依据意图即时生成：视觉学习者看到可操作图形，听觉学习者得到音频。它要求模型同时推断任务、用户偏好与可用设备；错误推断必须容易撤销，否则“隐形软件”也会让系统行为难以预测。

\lecturefigure{slide-58-dynamic-generative-ui.jpg}{Dynamic generative UI：界面随意图生成，软件创建趋于不可见}{00:47:40--00:48:04}

\lecturefigure{slide-59-personalized-health-education.jpg}{个性化医疗与教育加 consumer hardware 的未来入口}{00:48:04--00:48:42}

医疗与教育的价值很大，风险也不同。教育系统可以允许探索与纠错；医疗建议涉及症状遗漏、紧急分诊、隐私和监管。统一“个性化”指标会掩盖风险层级。产品必须区分 information、decision support 与 autonomous action，并为高风险输出设置专业复核和紧急升级。

\lecturefigure{slide-60-storytelling-relationship.jpg}{人与 storytelling process 的关系将因协作模型而改变}{00:48:42--00:49:18}

故事创作的开放性使它成为检验 social intelligence、voice preservation 与 agency 的好环境。模型可以给灵感、改写和结构建议，但“替作者决定什么值得表达”与“帮助作者实现意图”是不同产品。真正的 collaborator 应让作者看见来源、选择和改动，而不是用流畅文本夺走创作控制。

\begin{warningbox}{Future vignette 不等于 deployment claim}
这四张图表达讲者的方向判断。它们没有证明医疗安全、教育增益、动态 UI 可预测性或创作劳动影响已经解决。讲义把它们作为 research agenda：每个愿景都必须补上目标用户、失败模式、权限、验证、成本和责任人。
\end{warningbox}

\teachervoice{00:47:40--00:49:18，Karina 设想视觉/听觉偏好驱动的动态界面、个性化教育医疗和新的 storytelling 关系，同时希望创作者把 AI 当工具而非只感到威胁。课堂提示：这种乐观愿景必须与用户 agency 和验证机制一起实现。}

\subsection{Q\&A I：主观任务、benchmark 与创业路径}

Q\&A 首先追问：是否存在难以评测、因而研究者不愿碰的任务？讲者以 creative writing 和 emotional intelligence 为例，认为可以构造 benchmark，但复杂任务会更慢、更依赖里程碑。她也指出创业者不必先训练 foundation model；可以使用已有低成本模型，把创新集中在真实任务和产品闭环。

\begin{knowledgebox}{课堂提示：主观不等于不可评测}
00:49:18--00:55:00，Karina 区分“没有现成 frontier benchmark”和“原则上无法构造评测”。可从成对比较、专家 panel、用户长期 outcome、竞赛结果和反事实样例开始；但这些信号都携带价值判断，不能伪装成自然真值。
\end{knowledgebox}

\subsection{Q\&A II：如何避免 preference 收敛到全球平均}

讲者给出两条思路：保留 base model 的 entropy，以及用 RLAIF/synthetic preference data 定向生成需要的多样性。这里的 entropy 不是鼓励随机，而是避免 RL 把大量可用风格压成“平均用户喜欢的 emoji、Markdown 和语气”。

对上下文 $x$，策略 entropy 为
\[
\mathcal{H}(\pi(\cdot\mid x))=-\sum_y \pi(y\mid x)\log \pi(y\mid x).
\]
过强 preference optimization 可能降低 $\mathcal{H}$ 并造成 mode collapse。实践中应按任务保持候选多样性，同时用约束过滤事实错误和安全风险；“多样”不是“任何输出都好”。

\begin{importantbox}{课堂提示：RLAIF 的价值是可控覆盖，不是无限数据量}
00:55:00--01:01:20，Karina 说 synthetic data 不一定需要很多，关键是 diversity；团队仍可人工检查，或用校准较好的独立模型和 meta-eval 验证。她把 synthetic generation 描述为对目标分布的 curation，而不是复制平均人类偏好。
\end{importantbox}

\subsection{Q\&A III：定性诊断、成本与 frontier work}

对于模型 bug，讲者承认大量 nuanced weirdness 来自人工反复“玩模型”；自动 eval 适合检查已知 behavior，但新问题往往先由定性探索发现。她也区分产品开发与 frontier research 的成本：应用团队可以复用便宜模型，前沿团队必须承担发明阶段的低效率，之后才出现第二轮成本优化。

\begin{knowledgebox}{课堂提示：一次异常与系统性行为}
00:57:20--01:00:00，Karina 说定性检查要关注 behavior 是否持续出现。偶发 sample 可能来自随机性；跨 prompt 变体、会话和版本稳定复现，才应升级为系统性 issue。讲义建议把每个手工发现扩成变体集后进入自动回归。
\end{knowledgebox}

\begin{warningbox}{成本讨论的边界}
Q\&A 中关于训练成本、scaling 与未来降本的回答带有明显不确定性。可以确认的是：应用层可复用现有模型、前沿发明昂贵、后续工程会降本；不能从课堂回答推出训练成本必然按某条曲线下降。
\end{warningbox}

\subsection{Q\&A IV：机器人、AI coworker 与 social intelligence}

讲者对 robotics 保持乐观，但把数据称为巨大 bottleneck。对于“现在是否已有 AI 同事”，她明确回答尚未达到；理想形态更像共享屏幕、实时对话、指出具体对象并根据用户意愿调节 agency。缺口不是再加一个 Slack bot，而是 social intelligence、speech-to-speech、多模态 grounding 和共同工作空间。

\begin{knowledgebox}{课堂提示：AI coworker 还缺什么}
01:05:00--01:10:00，Karina 列出的缺口包括：实时生成与共同编辑、知道何时指导而非夺权、speech-to-speech、同时指向所谈对象，以及适配不同关系。这个清单比“能自动做任务”更严格，也解释了为什么 teammate eval 必须包含人际协调。
\end{knowledgebox}

\subsection{Q\&A V：Research-driven product 与传统 PRD 流程}

传统产品开发常从 PRD、设计和工程排期开始；research-driven product 可能先出现一个惊人的 capability demo，再围绕它塑造产品，也可能由产品和研究从第一天共同试验。Canvas 被讲者作为后者案例：post-training 与产品一起工作，过程更 ad hoc，但并不意味着没有规范；规范转移到了 prototype、eval、trace 和迭代证据。

\begin{importantbox}{课堂提示：能力 demo 不是产品终点}
01:09:52--01:11:29，Karina 对比传统软件流程和 research-driven product：研究先出现能力时，产品围绕能力成形；有时产品与研究从最初就共同设计。工程判断是：越 ad hoc，越需要记录 hypothesis、版本和 failure，避免团队只记得成功 demo。
\end{importantbox}

\subsection{本章小结}

未来愿景的共同条件不是“模型更聪明”，而是可信成功成本下降、界面能验证与撤销、个性化尊重用户控制、主观任务保留多样性、AI coworker 具备社会与多模态能力。Q\&A 让课程主张更完整：很多能力可教，但数据、验证、偏好分布和协作关系决定它是否成为可靠产品。

